\section{Chinchilla 与 LLaMA：先问清楚优化目标}

团队决定训练自有 LLM 时，Chinchilla 提供了当时最清晰的公开 scaling recipe。论文的经典结论是：在固定训练 compute 下，模型参数量和训练 token 应更均衡地扩展；其代表模型约为 70B 参数、1.4T tokens。课堂真正要纠正的不是这条结论，而是把它误用到另一个目标：\textbf{如果模型要被反复部署，是否应牺牲一次性训练成本来降低长期推理成本？}

\subsection{compute-optimal 与 inference-aware 的区别}

粗略地说，dense Transformer 预训练计算量可写成
\[
C_{\mathrm{train}}\approx 6ND,
\]
其中 $N$ 是参数量，$D$ 是训练 token 数。固定 $C_{\mathrm{train}}$ 时，选择 $N$ 与 $D$ 是 Chinchilla 问题；但一个产品的生命周期目标更接近
\[
J(N,D)=C_{\mathrm{train}}(N,D)
+Q\cdot C_{\mathrm{infer}}(N),
\]
其中 $Q$ 是部署期间累计生成量。若 $Q$ 很大，较小模型即使训练更多 tokens，也可能显著降低总成本和延迟。

\lecturefigure{05-objective-mismatch.png}{“Chinchilla trap”源于把固定训练预算目标误当成完整产品目标。}{本地字幕 09:20--11:40；Hoffmann et al. 2022；Touvron et al. 2023。}

\begin{knowledgebox}{读图：先确认预算约束，再讨论“最优”}
左侧最优意味着“这一次训练预算不能增加”；右侧最优意味着“训练只发生一次，推理会重复发生”。二者都可以合理，但不能混用。课程中所谓 over-train 小模型，是相对于固定 compute 的最优 token 数继续训练，而不是无上限重复同一数据。
\end{knowledgebox}

LLaMA 论文报告：6.7B 与 13B 模型训练约 1.0T tokens，32.5B 与 65.2B 模型训练约 1.4T tokens。这个设计强调，小模型在足够数据下可以获得远超早期直觉的能力，并在推理时更容易部署。它没有证明“参数越少越好”，而是把模型选择从单次 benchmark 转向质量、延迟、显存、吞吐和调用规模的联合问题。

\begin{warningbox}{“更多 tokens”不自动等于更强模型}
增加 $D$ 的收益取决于数据质量、重复率、学习率调度和模型容量。若新增 token 主要是低质量重复数据，训练可能只增加成本；若模型容量不足，边际收益也会下降。inference-aware training 仍需要 scaling experiment，而不是一句“把小模型训久一点”。
\end{warningbox}

\teachervoice{老师强调，许多人把“compute-optimal”当成普遍褒义词，却没有说明是在优化训练 loss、最终质量、推理延迟还是总拥有成本。工程讨论中遇到“optimal”，第一反应应是追问 objective、constraint 和 workload。}

\subsection{本章小结}

Chinchilla 给出固定训练预算下的配比规律；LLaMA 的课堂启示是把推理效率纳入目标。所谓 Chinchilla trap，本质是目标函数错配，而不是 scaling law 本身错误。

